\begin{longtable}{@{}L{1.05cm} L{6.0cm} L{7.2cm}@{}}
\caption{Điều nhóm rút ra được từ từng công trình và điều đó đổi gì trong thiết kế của nhóm.}\label{tab:baihoc}\\
\toprule \textbf{Mã} & \textbf{So với mô hình của nhóm} & \textbf{Bài học áp dụng} \\
\midrule\endfirsthead
\multicolumn{3}{@{}l}{\footnotesize\itshape Bảng \thetable{} (tiếp theo)}\\
\toprule \textbf{Mã} & \textbf{So với mô hình của nhóm} & \textbf{Bài học áp dụng} \\
\midrule\endhead
\midrule\multicolumn{3}{r@{}}{\footnotesize\itshape tiếp trang sau}\\\endfoot
\bottomrule\endlastfoot
\multicolumn{3}{@{}l}{\cellcolor{rowbg}\textbf{Học sâu dò fQRS}}\\[1pt]
p01 & {\scriptsize Không so sánh trực tiếp được ở mức con số vì đơn vị phân tích khác hẳn: Zhong tính F1 trên cửa sổ 100 ms đã cân bằng lớp, còn FetalQRS-TCN tính F1 trên nhịp tim với khớp ±50 ms; dù vậy bài vẫn là mốc đối chiếu quan trọng theo ba hướng. Về kiến trúc, Zhong dùng CNN 1D ba lớp với bộ lọc dài 3 nên tổng trường tiếp nhận chỉ vài chục mili giây trên đầu vào 100 ms, trong khi mô hình của nhóm có trường tiếp nhận 1.516 ms và khảo…} & {\scriptsize Áp dụng: đưa Zhong 2018 vào bài báo của nhóm như bằng chứng nguồn gốc cho luận điểm "tiền xử lý quan trọng hơn kiến trúc" (họ dùng 0,5–100 Hz, dải mà khảo sát của nhóm cho thấy mất 18,89 điểm F1 so với 10–60 Hz), học cách họ liệt kê rõ bản ghi của từng tập (a01–a07 kiểm thử, a08–a13 validation, 55 bản ghi còn lại huấn luyện) để nhóm cũng liệt kê bản ghi của từng fold tách theo bản ghi, và coi ý tưởng dùng SampEn làm chỉ số…} \\[2pt]
p09 & {\scriptsize Không so sánh trực tiếp được: họ có dùng ADFECGDB nhưng chỉ để tính AUC mức IMF (0,9282) và không bao giờ báo cáo F1 dò đỉnh R thai trên bộ này, nên không có con số nào đối chiếu được với macro F1 97,43 (tách theo bản ghi, 4 kênh) hay 99,18 (kênh tốt nhất) của mô hình của nhóm; hơn nữa họ không nêu dung sai ms nên dù có số cũng không biết họ đo ở ±50 ms hay rộng hơn. Về quy mô đánh giá và chi phí tính toán thì mô hình của…} & {\scriptsize Nên trích dẫn bài này làm bằng chứng độc lập cho luận điểm "tiền xử lý quan trọng hơn kiến trúc": họ khởi đầu bằng dải thông rộng 0,1-100 Hz nhưng đầu ra cuối cùng lại là Butterworth bậc 4 zero-phase 10-60 Hz, đúng dải mà nhóm chọn sau khi quét 8 dải (10-60 Hz thắng với 91,06), và họ cũng ghi nhận năng lượng QRS thai tập trung khoảng 8-45 Hz - khác biệt là một bài Sensors 2026 tìm ra dải này bằng tay còn nhóm chứng minh định…} \\[2pt]
p10 & {\scriptsize Không so sánh trực tiếp được, và mô hình của nhóm mạnh hơn rõ rệt ở mọi trục quan trọng: đơn vị phân tích khác hẳn (họ đánh giá trên ô 100 ms có tính cả ô âm, nhóm đánh giá trên nhịp với dung sai ±50 ms, nên 96,79\% accuracy của họ và macro F1 97,43 của nhóm là hai đại lượng khác loại, đặt cạnh nhau là nguy hiểm), họ dùng 4 kênh đồng thời còn nhóm dùng 1 kênh nên bài toán của nhóm khó hơn hẳn, họ chia ngẫu nhiên có xáo trộn…} & {\scriptsize Bài này là ví dụ để nhóm viết một đoạn về "tại sao giao thức đánh giá quan trọng hơn con số" trong phần Related Work hoặc Discussion, nêu đích danh rằng một bài Scientific Reports 2025 báo cáo 98,36\% accuracy nhưng đầu vào là 4 kênh chứ không phải 1 kênh, cửa sổ chồng lấp 90\% rồi xáo trộn trước khi chia, và đơn vị là ô 100 ms chứ không phải nhịp; đồng thời nên chỉ rõ rằng dùng khung 100 ms để xấp xỉ dung sai ±50 ms là một…} \\[2pt]
p12 & {\scriptsize Không so sánh trực tiếp được vì bốn lý do độc lập: họ dùng 12 kênh đồng thời còn nhóm dùng 1 kênh; họ phân loại nhị phân đoạn 65 ms còn nhóm định vị đỉnh với dung sai ±50 ms nên bài toán của nhóm khó hơn; nhãn chuẩn của họ là đầu ra thuật toán BSSR còn nhóm dùng ECG da đầu của ADFECGDB là tiêu chuẩn vàng thực sự; metric của họ là accuracy trên tập mất cân bằng chứ không phải macro F1 của lớp đỉnh. Nếu buộc phải đặt cạnh…} & {\scriptsize Nên sao chép cách phân tầng kết quả của họ (theo tuổi thai, BMI mẹ, có hay không biến chứng mẹ), bổ sung phân tích Bland-Altman kèm tương quan Pearson cho nhịp tim thai vì nhóm sẽ dễ dàng vượt xa r = 0,56, thêm hậu xử lý lọc nhịp kiểu SD-ROM lên chuỗi RR, và làm nổi bật việc nhóm báo cáo 113.481 tham số, 0,48 MB, 4,35 ms trong khi bài Q1 năm 2024 này không báo cáo gì về chi phí tính toán. Phải tránh việc dùng accuracy làm…} \\[2pt]
\multicolumn{3}{@{}l}{\cellcolor{rowbg}\textbf{Tách nguồn / tái tạo dạng sóng}}\\[1pt]
p03 & {\scriptsize Trên cùng bộ ADFECGDB và cùng kiểu giao thức bỏ một đối tượng, họ đạt F1 99,7 ± 0,4\% ở dung sai 30 ms, còn nhóm đạt 99,18\% (Se 99,59, PPV 98,79) trên kênh lâm sàng tốt nhất và 97,43 ± 4,37\% trên toàn bộ 4 kênh ở dung sai 50 ms; vì mỗi bản ghi ADFECGDB là một sản phụ nên tách theo bản ghi của nhóm tương đương tách theo chủ thể của họ, do đó họ nhỉnh hơn nhóm (khoảng 0,5 điểm trên kênh tốt nhất, khoảng 2,3 điểm trên trung bình…} & {\scriptsize Áp dụng ngay: bổ sung báo cáo kết quả ở dung sai 30 ms bên cạnh 50 ms cho tất cả thực nghiệm để đặt cạnh trực tiếp con số 99,7\% và 94,7\% của họ (nếu chỉ báo cáo 50 ms, phản biện Q1 sẽ nói nhóm chọn dung sai để đẹp số), bổ sung khoảng tin cậy 95\% cho macro F1 theo chuẩn STARD/TRIPOD thay vì chỉ có trung bình ± độ lệch chuẩn, đưa con số 94,7\% xuyên bộ dữ liệu của họ vào bảng so sánh và giải thích thẳng thắn vì sao 74,28 \% (PSD mù nhãn, 60 bản ghi sạch; các con số 79,40 trên 75 bản ô
nhiễm và 77,34 trên mẫu 10 bản ghi \textbf{đều đã bị rút}) của…} \\[2pt]
p04 & {\scriptsize Không so sánh trực tiếp được vì ba lý do độc lập mà mỗi lý do đều đủ phá vỡ phép so sánh: họ dùng 4 kênh bụng làm đầu vào bộ sinh còn FetalQRS-TCN dùng đơn kênh, nên 96,4\% của họ là kết quả của một bài toán dễ hơn hẳn; giao thức 5-fold không nêu cách nhóm và bằng chứng số học ở Bảng 5 (nhịp thật giống hệt nhau ở cả 5 fold) chỉ rất mạnh tới việc chia ngẫu nhiên theo đoạn trên dữ liệu gộp, tức là có rò rỉ dữ liệu, trong khi…} & {\scriptsize Áp dụng: dùng con số 0,762 triệu tham số của họ làm mốc đối chiếu chính cho luận điểm hiệu quả tính toán của nhóm (FetalQRS-TCN chỉ bằng 1/3 một bộ sinh của họ và 1/6,7 toàn mạng mà vẫn không cần fECG da đầu làm nhãn), học cách họ báo cáo chi phí tính toán với phần cứng cụ thể (Xeon 4 nhân 8 luồng 2,0 GHz, 25 GB RAM, Tesla T4 16 GB), số tham số tách riêng từng thành phần và thông lượng (1.740 giây tín hiệu / 73 giây) rồi bổ…} \\[2pt]
p05 & {\scriptsize Không so sánh trực tiếp được hoàn toàn, nhưng đây là công trình gần nhất với giao thức của nhóm ta và là mốc so sánh đáng tin nhất trong ba bài: cùng dung sai 50 ms, cùng cửa sổ 4 giây (họ tìm ra 4 giây là tối ưu bằng thực nghiệm, trùng khớp với kết quả quét trường tiếp nhận của nhóm ta là bão hoà quanh 1,5 s), cùng định hướng một kênh và cùng cách chọn kênh tốt nhất. Hai con số không thể đặt cạnh nhau vì khác tập dữ liệu và…} & {\scriptsize Nên chạy ngay thí nghiệm ưu tiên số một là huấn luyện chính FetalQRS-TCN trên FECGSYNDB (145,8 giờ, 250 Hz, đúng tần số nhóm ta đang dùng nên không cần lấy mẫu lại) rồi đo trên ADFECGDB và CinC 2013 mà không tinh chỉnh, kèm ba kiểu tăng cường dữ liệu của họ (jittering với công suất nhiễu thấp hơn năng lượng QRS thai 3 dB, cắt ngẫu nhiên theo phân phối Beta, che thời gian–tần số chỉ trên các ô dưới 50 Hz) vốn mang lại +4,35…} \\[2pt]
p06 & {\scriptsize Không so sánh trực tiếp được, vì con số 99,17\% trên ADFECGDB của họ được tạo ra sau khi tinh chỉnh trên chính bộ ADFECGDB, còn 97,43 ± 4,37 (toàn bộ 4 kênh) và 99,18 (kênh tốt nhất) của nhóm ta là kết quả tách theo bản ghi, trong đó mọi bản ghi kiểm thử đều chưa từng xuất hiện khi huấn luyện; đặt hai con số cạnh nhau là so một kết quả có thể có rò rỉ với một kết quả sạch, về nguyên tắc không được làm. Thêm ba khác biệt: dung…} & {\scriptsize Nên học ba điểm: báo cáo kết quả theo từng bản ghi như Bảng 2 và Bảng 5 của họ (R01, R04, R07, R08, R10 riêng từng dòng) thay vì chỉ báo cáo macro F1 97,43 ± 4,37 vì độ lệch chuẩn 4,37 đang giấu đi bản ghi yếu; báo cáo F1 ở cả ba mức dung sai — 50 ms (chuẩn CinC, để so sánh), 31,25 ms (để đối chiếu với bài này) và 150 ms (AAMI EC57) — thành một bảng độ nhạy theo dung sai, vì với jitter chỉ 3,05 ms của FetalQRS-TCN thì nhóm…} \\[2pt]
p07 & {\scriptsize Đây là mốc so sánh tốt nhất trong ba bài và đối chiếu được một phần, vì bốn tham số trùng khớp hoàn toàn với FetalQRS-TCN: dung sai 50 ms, đầu vào một kênh, cửa sổ khoảng 4 giây và tần số làm việc 250 Hz. Khác biệt nằm ở nguồn huấn luyện và đơn vị phân tích: họ huấn luyện 100\% trên dữ liệu tổng hợp rồi kiểm thử zero-shot còn nhóm ta huấn luyện tách theo bản ghi trong nội bộ ADFECGDB, nên 96,52\% của họ và 97,43 ± 4,37 của…} & {\scriptsize Thí nghiệm ưu tiên số một là huấn luyện FetalQRS-TCN trên dữ liệu tổng hợp FECGSYN theo đúng công thức của bài này — 55.000 đoạn chia theo SINR thành ba nhóm (20.000 chất lượng cao với SINR\_f trong [−5; 0], 25.000 trung bình với SINR\_f trong [−15; −5], 10.000 thấp với SINR\_f trong [−25; −15]), các kịch bản sinh lý và nhiễu MA/BW/EM, trong đó nhóm ta có thể dùng cả 7 kịch bản của FECGSYNDB thay vì 3 (thêm co thắt tử cung C3,…} \\[2pt]
p11 & {\scriptsize Đây là bài duy nhất trong ba bài so sánh trực tiếp được với mô hình của nhóm, và kết quả gần như ngang nhau: CUNet đạt F-score 77,8 ± 18,6 trên CinC 2013 zero-shot không tinh chỉnh, còn nhóm đạt 74,28 ± 30,57 (PSD mù nhãn) hoặc 82,01 (quy tắc \texttt{peakprob} hậu kiểm) trên \textbf{60 bản ghi sạch} của CinC 2013 set-a, xuyên bộ dữ liệu không tinh chỉnh (các con số 79,40 ± 29,18 trên 75 bản ô nhiễm và 77,34 ± 28,54 trên mẫu 10 bản ghi \textbf{đều đã bị rút}); lưu ý hai con số không nằm trên cùng tập vì Orvas chấm đủ 75 bản, trong đó 15 bản mô hình của nhóm đã thấy khi huấn luyện, với các điều kiện khớp nhau về dung sai ±50 ms, đơn vị phân tích là nhịp, công thức F-score = 2TP/(2TP+FN+FP), đầu vào…} & {\scriptsize Đây là mốc chuẩn chính xác nhất mà nhóm cần: hãy đưa con số 77,8 ± 18,6 của Orvas 2025 vào bảng so sánh của nhóm ngay cạnh 74,28 (PSD mù nhãn) và 82,01 (\texttt{peakprob} hậu kiểm) trên 60 bản ghi sạch, kèm ghi chú rõ ràng là cùng ±50 ms, cùng đơn kênh, cùng zero-shot trên CinC 2013 nhưng \textbf{không cùng tập bản ghi} - điều này đặt kết quả của nhóm ở mức "ngang bằng trình độ hiện tại năm 2025" mà không cần con số 79,40 đã rút, đồng thời nên bổ sung chỉ số HRerr (bpm) vì đó là sai số nhịp tim thai mà bác sĩ thực sự quan tâm và cho phép so…} \\[2pt]
\multicolumn{3}{@{}l}{\cellcolor{rowbg}\textbf{Xử lý tín hiệu cổ điển}}\\[1pt]
p16 & {\scriptsize Không so sánh được, và lý do rất quan trọng về mặt khái niệm vì đây không phải cùng một bài toán: FetalQRS-TCN là bộ dò đỉnh (đầu vào là tín hiệu, đầu ra là vị trí fQRS, đo bằng F1 với dung sai 50 ms) còn par-EKS là bộ tách/tăng cường tín hiệu cần sẵn vị trí đỉnh R thai làm đầu vào, nên mô hình của nhóm có thể dùng làm đầu vào cho par-EKS chứ không cạnh tranh với nó; nếu phản biện hỏi vì sao không so với Niknazar 2013 thì…} & {\scriptsize Dùng bài này làm ví dụ điển hình trong phần tổng quan về việc phần lớn công trình đơn kênh không báo cáo chỉ số phát hiện đỉnh chuẩn hóa, với cách viết bám sát văn bản gốc rằng Niknazar và cộng sự (2013) báo cáo mức cải thiện SNR trên 20 dB trên dữ liệu tổng hợp nhưng không báo cáo Se/PPV/F1 trên dữ liệu lâm sàng, khiến kết quả của họ không thể đối chiếu trực tiếp với các bộ dò đỉnh; đồng thời khai thác hướng mở rộng ghép…} \\[2pt]
p17 & {\scriptsize Đây là bài so sánh trực tiếp được nhất trong ba bài vì cùng bộ dữ liệu ADFECGDB, cùng dung sai 50 ms, cùng công thức F1 và thực sự đơn kênh: trên ADFECGDB với tất cả các kênh, Castillo đạt 94,11\% (17/20 tín hiệu, trong mẫu) so với 97,43 ± 4,37 của FetalQRS-TCN (tách theo bản ghi, 20 đánh giá trên đủ 4 kênh) nên nhóm cao hơn khoảng 3,3 điểm và giao thức của nhóm nghiêm ngặt hơn, còn ở kênh tốt nhất là 98,63\% so với 99,18 của…} & {\scriptsize Nên đưa bài này vào bảng so sánh chính với chú thích đầy đủ (Castillo 2018, ADFECGDB, đơn kênh, dung sai 50 ms, F1 98,04\% trên 11 tín hiệu do bác sĩ chọn và 94,11\% trên 17 tín hiệu, không tách theo bản ghi) vì đây là baseline cổ điển mạnh nhất và sát nhất với bài toán của nhóm, đồng thời bắt chước cách báo cáo ba tầng của họ (toàn bộ, nhóm sạch, kênh tốt nhất) kèm thêm cột tỉ lệ đoạn được chấp nhận, thêm bước sửa FP/FN theo…} \\[2pt]
p19 & {\scriptsize Không so sánh trực tiếp được, vì bốn lý do độc lập: họ dùng NInFEA riêng tư với 20 thai phụ tuần 21-27 và nhãn chuẩn từ Doppler xung còn nhóm dùng ADFECGDB với 5 sản phụ chuyển dạ và nhãn chuẩn từ điện cực da đầu thai nhi; độ đo của họ là Acc = (TP+TN)/(TP+TN+FP+FN) không rõ cách đếm TN trong khi nhóm báo cáo macro F1, Se và PPV; họ không nêu dung sai khớp nhịp nên không biết khoan dung hơn hay chặt hơn mức ±50 ms của nhóm;…} & {\scriptsize Trích dẫn bài này làm mốc đối chứng cho luận điểm "1 kênh cộng học sâu" trong phần Related Work (lọc thích nghi đơn tham chiếu chỉ đạt khoảng 68\% độ chính xác trên dữ liệu thật dù đã thêm điện cực ngực, trong khi nhóm đạt 97,43\% F1 mà không cần điện cực ngực nào), dùng bề rộng QRS thai 40 ms và QRS mẹ 100 ms mà họ dẫn từ Taylor 2003 và Surawicz 2009 để biện minh cho sigma = 12 ms của nhãn Gaussian (khoảng ±2 sigma bằng 48…} \\[2pt]
\multicolumn{3}{@{}l}{\cellcolor{rowbg}\textbf{Đối chuẩn, dữ liệu, thử thách}}\\[1pt]
p13 & {\scriptsize Điểm chung quan trọng nhất là cùng dung sai ±50 ms và cùng công thức F1 = 2TP/(2TP+FN+FP), nên con số của Behar và của nhóm nằm trên cùng một thang đo — điều mà rất ít bài trong danh sách đạt được. Tuy vậy chỉ so sánh được một phần: Chương 6 dùng 1 kênh bụng nhưng cộng thêm 1 kênh ngực mẹ còn Chương 7 dùng 4 kênh bụng nên đều nhiều thông tin hơn nhóm; Behar không bao giờ báo cáo F1 riêng cho ADFECGDB vì 25 bản ghi ADFECGDB…} & {\scriptsize Việc phải làm ngay là kiểm tra chồng lấn giữa ADFECGDB (r01, r04, r07, r08, r10) và CinC 2013 set-a bằng tương quan chéo tín hiệu hoặc chuỗi RR tham chiếu trước khi công bố con số CinC nào, vì nếu không kiểm tra thì reviewer Q1 có thể tự phát hiện và bài sẽ bị từ chối vì rò rỉ dữ liệu. \textbf{[Sửa vòng 7: đã kiểm --- 15/75 bản set-a là bản sao ADFECGDB; các con số 79,40 (75 bản) và 77,34 (mẫu 10 bản) đều đã rút; số hiện hành là 74,28 trên 60 bản sạch. Ghi chú p13 này và ghi chú p20 đều đã nêu đúng việc cần làm, và nhóm đã bỏ sót.]} Nên trích dẫn Behar làm căn cứ cho dung sai ±50 ms (AAMI EC57 cho phép 150 ms nhưng đó là ngưỡng cho ECG người lớn, còn nhịp thai nhanh gấp…} \\[2pt]
p14 & {\scriptsize Bài cùng thang đo với nhóm ở một điểm quan trọng là dung sai ±50 ms và công thức F1 giống hệt, nhưng chỉ so sánh được một phần vì Behar dùng 4 kênh bụng cộng 5 phương pháp tách chạy song song và cơ chế chọn kênh trong khi nhóm dùng 1 kênh và 1 mô hình duy nhất — đó là lợi thế cấu trúc chứ không phải lợi thế thuật toán — và vì F1 = 96,0\% của ông là trong mẫu trên chính tập quét tham số còn 97,43 ± 4,37 tách theo bản ghi của…} & {\scriptsize Áp dụng ngay với chi phí thấp: hậu xử lý làm mượt RR của Behar (lấy trung vị RR của 5 nhịp trước, chỉ làm mượt khi trung vị nằm trong 110–170 bpm, xoá nhịp cách nhịp trước < 70\% trung vị RR và chèn nhịp khi > 175\% trung vị RR), cơ chế BCM để loại chuỗi nhịp mẹ bị nhận nhầm thành nhịp thai, notch có điều kiện cho dữ liệu trộn nguồn 50/60 Hz, và hiệu chỉnh đỉnh R mẹ theo từng kênh trong ±30 ms trước khi trừ mẫu trung vị. Khi…} \\[2pt]
p15 & {\scriptsize Không so sánh trực tiếp được về con số, và đây là điều phải nói rõ trong bài của nhóm: họ đo trên dữ liệu mô phỏng còn nhóm đo trên ADFECGDB lâm sàng thật, F1 99,9\% của BSSica là cận trên oracle sau khi chọn kênh tốt nhất theo nhãn chuẩn trong 8 kênh, và BSS cần 8 kênh còn AM cần thêm kênh tham chiếu mẹ trong khi FetalQRS-TCN chỉ cần 1 kênh bụng. Tuy vậy vẫn có những trục so sánh được và có lợi cho nhóm: cùng dung sai 50 ms…} & {\scriptsize Nên tải FECGSYNDB (miễn phí trên PhysioNet, tần số 250 Hz trùng khớp với tần số đích của pipeline nên nạp thẳng không cần lấy mẫu lại) để chạy FetalQRS-TCN trên 7 kịch bản phi dừng × 5 mức SNR, luôn báo cáo F1 kèm MAE/jitter thành cặp đúng như tác giả khuyến nghị, báo cáo tỉ lệ đoạn bị loại theo tiền lệ 69,7\% xuống 11,6\% của họ, và trích dẫn bài này vừa để biện minh cho dung sai 50 ms vừa làm căn cứ cho ý tưởng fSQI (họ viết…} \\[2pt]
p20 & {\scriptsize So sánh được một phần và theo hướng ngoài ý muốn: bộ dữ liệu thì so sánh được vì đây chính là tập nhóm dùng để đo xuyên bộ dữ liệu (74,28 F1 trên 60 bản ghi sạch; con số cũ 79,40 ± 29,18 trên toàn bộ 75 bản ghi đã rút), nhưng con số thì không so sánh được vì tiêu chí chính thức của CinC 2013 là lỗi FHR tính bằng (nhịp/phút)\textasciicircum{}2 và lỗi RR tính bằng ms chứ không phải F1 với dung sai 50 ms. Điểm phải sửa ngay trong bài của nhóm là giả định "50 ms là chuẩn CinC 2013" không được bài xã luận gốc này…} & {\scriptsize Việc đầu tiên là kiểm tra chồng lấn giữa ADFECGDB và CinC 2013 bằng cách đối chiếu từng bản ghi trong tập đang test với 25 bản ghi ADFECGDB có trong cuộc thi, loại bỏ mọi bản ghi trùng rồi tính lại con số CinC của nhóm, và nếu không loại được thì phải ghi rõ cảnh báo trong bài kèm trích dẫn chính câu của Clifford về trường hợp Rodrigues; đồng thời sửa cách trích dẫn dung sai 50 ms sang đúng nguồn (benchmark của Behar) hoặc… \textbf{[Sửa vòng 7: đã kiểm --- 15/75 bản set-a là bản sao ADFECGDB; 79,40 đã rút, số đúng 74,28 trên 60 bản sạch; chính câu cảnh báo của Clifford ở đây là điều nhóm đã bỏ sót.]}} \\[2pt]
p21 & {\scriptsize So sánh được một phần và đây là bài quan trọng nhất trong ba bài đối với nhóm, vì bộ B2\_Labour chứa chính 5 bản ghi ADFECGDB mà nhóm đang dùng (r01, r04, r07, r08, r10), cùng điều kiện chuyển dạ và cùng nhãn chuẩn là FECG trực tiếp từ đầu thai nhi: con số gần nhất để đối chiếu là PCA đạt F1 = 98,56\% và ICA đạt F1 = 98,55\% trên toàn bộ 12 bản ghi B2 với dung sai ±40 ms, trong khi mô hình của nhóm đạt macro F1 = 97,43 ± 4,37…} & {\scriptsize Ưu tiên cao nhất là tải bộ B2 đầy đủ 12 bản ghi từ figshare (DOI 10.6084/m9.figshare.c.4740794, giấy phép CC BY 4.0) để lấy 7 bản ghi chuyển dạ chưa từng nằm trong ADFECGDB làm tập kiểm tra ngoài lý tưởng (cùng phân bố lâm sàng nên loại bỏ được yếu tố khác bộ dữ liệu, nhưng mô hình chưa từng thấy), và tải bộ B1 gồm 10 bản ghi thai kỳ 20 phút tuần 32-42 để mở rộng phạm vi ứng dụng từ chuyển dạ sang trước sinh, với điều kiện…} \\[2pt]
\multicolumn{3}{@{}l}{\cellcolor{rowbg}\textbf{Lý thuyết TDA}}\\[1pt]
p22 & {\scriptsize Không so sánh trực tiếp được vì ba lý do: bài toán khác hẳn (họ chấm điểm tuần hoàn toàn cục cho cả đoạn tín hiệu, còn nhóm làm định vị đỉnh từng mẫu theo kiểu sequence-to-sequence), dữ liệu khác hẳn (tổng hợp 50 mẫu vô thứ nguyên, không phải ADFECGDB hay CinC 2013), và không có Se/PPV/F1 theo nhịp, không có dung sai ±50 ms, không có giao thức tách theo bản ghi. Tuy nhiên đây là cơ sở lý thuyết mạnh nhất để bảo vệ phát hiện…} & {\scriptsize Nên áp dụng: nếu dùng TDA làm chỉ số chất lượng tín hiệu thai (fSQI) thì phải đặt kích thước cửa sổ trượt xấp xỉ một chu kỳ RR thai (khoảng 400-500 ms, tương ứng khoảng 100-125 mẫu ở 250 Hz) chứ không được ngắn hơn — căn cứ là Corollary 6.10 và Prop 5.5 — đồng thời chọn số chiều nhúng M ≥ 2N với N là bậc hài cao nhất muốn giữ, chọn số nguyên tố p > N khi tính đồng điều bền vững để tránh hiện tượng xoắn kiểu Möbius, bắt buộc…} \\[2pt]
p23 & {\scriptsize Không so sánh trực tiếp được với FetalQRS-TCN vì miền dữ liệu hoàn toàn khác (đám mây điểm tổng hợp và mặt PDE, không phải ECG bụng), không có Se/PPV/F1 theo nhịp, không có dung sai ms, không có giao thức tách theo bản ghi hay tách theo chủ thể. Tuy nhiên đây là bài báo gốc của biểu diễn mà nhóm đã thử nghiệm và chốt bảng (cnn2d\_PersistenceImage đạt 27,42 so với 97,14 của CNN 1D cùng ngân sách tham số), và điểm mấu chốt để…} & {\scriptsize Nên áp dụng: nếu vẫn dùng PI cho fSQI thì phải đóng khung lại thành bài toán một nhãn cho một cửa sổ đúng thiết kế gốc của PI chứ không dùng cho sequence-to-sequence, đồng thời nối PI của H0 và H1 lại với nhau vì bài chứng minh lợi ích lớn (82,5\% so với 49,8\% khi chỉ H0 và 65,7\% khi chỉ H1 trên linked twist map; 97,3\% so với 94,7\% và 93,3\% trên aKS), không cần quét siêu tham số lớn mà có thể cố định 20×20 và σ = 0,1 do độ…} \\[2pt]
p24 & {\scriptsize Không so sánh trực tiếp được với FetalQRS-TCN vì miền dữ liệu hoàn toàn khác (đồ thị mạng xã hội, phân tử hóa học, quỹ đạo hệ động lực — không có tín hiệu sinh lý), không có Se/PPV/F1 theo nhịp, không có dung sai ±50 ms, không có tách theo bản ghi hay tách theo chủ thể, và không có số tham số để đối chiếu với 113.481 tham số của nhóm. Ý nghĩa gián tiếp lại rất lớn cho lập luận của nhóm: PersLay là nỗ lực mạnh nhất hiện nay…} & {\scriptsize Nên áp dụng: nếu nhóm quay lại TDA cho fSQI thì PersLay là lựa chọn đúng đắn hơn Persistence Image cố định vì lưới trọng số học được còn cho phép diễn giải vùng nào của sơ đồ quan trọng sau huấn luyện (Hình 5 phụ lục) — rất có giá trị cho bài báo Q1 — đồng thời nên xét dùng bền vững mở rộng thay vì bền vững thường vì nó bắt được cả nhánh hướng lên chứ không chỉ đỉnh và mọi điểm có tọa độ hữu hạn (chênh lệch 14,9 điểm trên…} \\[2pt]
p25 & {\scriptsize Không so sánh trực tiếp được vì khác bài toán (nhận dạng hoạt động và phân loại ảnh chứ không phải dò đỉnh fQRS), khác dữ liệu (gia tốc kế đeo tay và ảnh CIFAR/SVHN, không có ECG bụng), khác chỉ số (F1 có trọng số đa lớp so với macro F1 khớp nhịp ±50 ms) và không có khái niệm dung sai ms. Tuy vậy bài củng cố gián tiếp kết quả của nhóm: ngay trong một công trình chủ trương dùng TDA, PI dùng một mình (MLP-PI 46,45) vẫn thua xa…} & {\scriptsize Nếu nhóm vẫn giữ một nhánh TDA cho fSQI thì không nên tính PI bằng thư viện giải tích trong vòng lặp huấn luyện mà tính trực tiếp peak prominence (tương đương H0 lọc theo tập mức), đồng thời có thể mượn cách họ giảm chiều PI bằng PCA xuống 32 chiều rồi nối vào đầu ra global-average-pooling để đưa điểm fSQI vào FetalQRS-TCN mà không làm phồng số tham số. Bảng thời gian của họ (3567 ms cho một ảnh 32 x 32) là bằng chứng định…} \\[2pt]
p26 & {\scriptsize Không so sánh trực tiếp được về con số vì khác hoàn toàn bài toán (phân loại hình dạng từ đám mây điểm, không có ECG, không có đỉnh, không có dung sai ms, không có bệnh nhân). Nhưng đây là căn cứ quan trọng nhất để biện minh cho quyết định chuyển hướng của nhóm: bộ hướng dẫn của họ liệt kê đúng ba loại tín hiệu mà PH làm tốt — số chu trình k chiều (lọc alpha, khoảng dài nhất), độ cong (Vietoris-Rips, nhiều khoảng ngắn) và…} & {\scriptsize Trích bộ hướng dẫn và phần kết luận của Turkes và cs. làm căn cứ học thuật chính thức cho luận điểm "TDA không phù hợp với bài toán dò vị trí đỉnh", đồng thời áp dụng đúng khung ba bước đó cho nhánh fSQI: xác định tín hiệu (có tồn tại vòng lặp tuần hoàn ứng với nhịp tim thai hay không), chọn phép lọc (nhúng trễ để biến tính chu kỳ thành vòng lặp H1), chọn đặc trưng (max persistence hoặc tổng life-time của H1), và dùng lọc…} \\[2pt]
\multicolumn{3}{@{}l}{\cellcolor{rowbg}\textbf{Nhúng trễ và chọn tham số}}\\[1pt]
p27 & {\scriptsize Chỉ so sánh được một phần: họ có dùng hai bộ ECG thai nhi (NonInvasiveFetalECGThorax1 và Thorax2, mỗi bộ 3.765 chuỗi dài 750 mẫu) và đạt 91,8\% và 93,0\%, nhưng đó là phân loại 42 lớp hình thái do chuyên gia gán nhãn trên toàn bộ đoạn chuỗi chứ không phải dò vị trí phức bộ QRS thai nhi, không có dung sai ms, không có Se/PPV theo nhịp, không có jitter và không có khái niệm khử ECG mẹ, nên 91,8\% của họ và 99,18\% macro F1 của…} & {\scriptsize Áp dụng trực tiếp cho fSQI: thay vì chọn một độ trễ, hãy quét tau qua một dải rồi xếp chồng thành PD ba chiều — với ECG bụng 250 Hz và nhịp tim thai 120–160 lần/phút (chu kỳ khoảng 375–500 ms, tương ứng 94–125 mẫu), dải tau nên phủ từ khoảng 1/10 đến 1/2 chu kỳ này, cửa sổ phải bao ít nhất 2 chu kỳ (khớp với cửa sổ 4 giây đang dùng), và nên dùng kết hợp hạt nhân K = alpha\_0·K0 + alpha\_1·K1 với alpha\_0 chọn bằng kiểm định…} \\[2pt]
p28 & {\scriptsize Không so sánh được về số liệu với mô hình của nhóm vì bài không làm dò đỉnh, không có ECG, không có F1/Se/PPV, không có dung sai theo ms và không có chủ thể để tách theo bản ghi, nhưng lại rất có giá trị làm cơ sở lý thuyết cho quyết định chuyển hướng. Bài chứng minh bằng chính cấu trúc thuật toán rằng muốn có lỗ H1 trong nhúng trễ 2 chiều thì đoạn tín hiệu phải dài ít nhất l = 4τ, tức khoảng một chu kỳ đầy đủ, và còn phải…} & {\scriptsize Nên trích dẫn ràng buộc độ dài strand l = 4τ và ràng buộc Nhole ≥ 8 làm căn cứ học thuật cho lập luận "cửa sổ ngắn hơn một chu kỳ tim thì không thể chứa vòng nhịp" khi giải thích vì sao Persistence Image kết hợp CNN 2D đứng chót bảng (27,42 so với 97,14), đồng thời dùng Bảng I về thời gian tính toán làm dẫn chứng số thật cho phần chi phí; theo heuristic một phần tư chu kỳ, chu kỳ tim thai khoảng 430 ms cho τ khoảng 107 ms,…} \\[2pt]
\multicolumn{3}{@{}l}{\cellcolor{rowbg}\textbf{TDA cho tín hiệu tim}}\\[1pt]
p29 & {\scriptsize So sánh được một phần và là bài tham chiếu quan trọng nhất cho hướng fSQI, nhưng không so được về con số với FetalQRS-TCN vì nhiệm vụ khác hẳn (họ phân loại chất lượng đoạn 10 giây, nhóm định vị từng đỉnh fQRS với dung sai ±50 ms), đối tượng khác (ECG người lớn, không có thai nhi và không có bài toán khử ECG mẹ), chỉ số khác (F1 phân loại nhị phân trên đoạn so với F1 dò nhịp có dung sai thời gian), và họ có rò rỉ dữ liệu…} & {\scriptsize Nên trích dẫn bài này làm nền cho hướng fSQI vì nó cho thấy đồng điều bền vững có giá trị thật cho bài toán đánh giá chất lượng (mAcc 98,04 bằng SubLevel-Set) trong khi thất bại cho bài toán định vị đỉnh, đồng thời bắt chước cấu trúc bảng đối chiếu với các chỉ số chất lượng kinh điển (bSQI, MS-QI, picaSQI, kSQI, pSQI) và ghi nhận rằng SubLevel-Set — loại lọc nhóm đã chứng minh tương đương peak prominence — vẫn đủ mạnh (97,25…} \\[2pt]
p30 & {\scriptsize So sánh được một phần vì cùng lĩnh vực ECG và cùng dùng đồng điều bền vững lọc sublevel-set, nhưng không so được về con số: họ phân loại nhãn của một nhịp đã biết vị trí còn nhóm định vị vị trí đỉnh chưa biết, họ không dò đỉnh nên không có dung sai ±50 ms và không có Se hay PPV theo nghĩa dò nhịp, chỉ số của họ là độ chính xác có trọng số còn của nhóm là macro F1 dò nhịp (hai đại lượng không quy đổi được), và đối tượng là…} & {\scriptsize Nên trích dẫn định lý bất biến đường cong Betti làm lập luận trung tâm giải thích vì sao Persistence Image kết hợp CNN 2D đứng chót bảng 16 kiến trúc (27,42 so với 97,14), học cách trình bày điểm kiểm định tách biệt điểm kiểm tra theo bệnh nhân để phơi bày đúng độ lớn khác biệt cá thể (tương ứng với nhóm là 97,43 ± 4,37 khi tách theo bản ghi so với 74,28 ± 30,57 khi chạy xuyên bộ dữ liệu CinC 2013 trên 60 bản ghi sạch; cặp số cũ 79,40 ± 29,18 trên 75 bản ghi đã rút), bắt chước ý tưởng kênh…} \\[2pt]
p31 & {\scriptsize Không thể so sánh trực tiếp về con số vì đây là bài toán khác hẳn — phân loại giai đoạn ngủ của người lớn từ chuỗi HRV đã có sẵn, đo bằng Acc/Kappa/AUC trên epoch 30 giây — trong khi nhóm dò đỉnh fQRS từ ECG bụng thai nhi và đo bằng F1 khớp nhịp với dung sai ±50 ms; nhưng nếu đọc về mức độ thành công thì F1 cao nhất của họ là 0,510 (REM/NREM) và 0,452 (Thức/Ngủ), so với 97,43 ± 4,37 macro F1 của FetalQRS-TCN, cho thấy TDA…} & {\scriptsize Nên áp dụng: trích dẫn bài này (Frontiers in Physiology 2021) làm tiền lệ cho giao thức kiểm định chéo cơ sở dữ liệu không tinh chỉnh mà nhóm đã chạy trên CinC 2013, nhấn mạnh trong phần Methods rằng chính họ tuyên bố trừ bài [50] ra thì không ai báo cáo cross-database; mượn nguyên 16 persistence statistics (trung bình, độ lệch chuẩn, độ lệch, độ nhọn, phân vị 25/50/75, entropy bền vững trên cả hai tập M và L) làm ứng viên…} \\[2pt]
p32 & {\scriptsize Không thể so sánh trực tiếp, lý do rất đơn giản: bài không có một con số hiệu năng nào để đặt cạnh macro F1 97,43 ± 4,37 của FetalQRS-TCN vì họ không dò đỉnh, không phân loại, không có tập kiểm tra, toàn bộ kết quả là p-value và hệ số tương quan mức nhóm; nhiệm vụ cũng khác hẳn — mô tả sự trưởng thành hệ thần kinh tự chủ ở trẻ em từ chuỗi RR sạch đã được chuyên gia sửa tay, trong khi nhóm làm việc với ECG bụng đơn kênh đầy…} & {\scriptsize Nên áp dụng: trích dẫn bài này (PLoS One tháng 12/2025, Universidad Politecnica de Madrid) làm khoảng trống nghiên cứu cho phần Giới thiệu vì chính họ tuyên bố mở rộng TDA sang HRV thai nhi là việc chưa ai làm — đúng chỗ đứng của nhóm, với dữ liệu ADFECGDB sẵn có; mượn cách trình bày minh bạch về tham số nhúng (báo cáo cả kết quả AMI với τ tối ưu 7-12 nhịp, FNN d ≈ 6, tiêu chí Cao d ≈ 9 rồi mới nói chọn gì và vì sao, dù lựa…} \\[2pt]
p33 & {\scriptsize Không thể so sánh trực tiếp về con số: đây là bài toán phân loại trạng thái căng thẳng trên cửa sổ 60 giây của người lớn, đo bằng độ chính xác và F1 ở mức cửa sổ, còn FetalQRS-TCN là bài toán dò đỉnh fQRS trên ECG bụng thai nhi đo bằng F1 khớp nhịp với dung sai ±50 ms, không có điểm chung nào về nhãn hay dung sai. Nhưng về giao thức thì đối chiếu được một phần và rất có giá trị: tách theo chủ thể của họ tương đương tách theo…} & {\scriptsize Nên áp dụng: trích dẫn cặp số 96\% (nội bộ đối tượng) so với 81\% (tách theo chủ thể), p < 0,05, WESAD ba lớp làm viện dẫn tốt nhất trong cả ba bài để biện minh vì sao nhóm bắt buộc dùng tách theo bản ghi và tách theo chủ thể, và vì sao con số 91\% trong README cũ là không hợp lệ — nên đưa thẳng vào phần Methods hoặc Discussion; mượn kỹ thuật subwindowing kèm công thức cập nhật tăng dần μ\_\{N+1\} = μ\_N + (f\_\{M+1\} − f\_1)/M cho…} \\[2pt]
\end{longtable}